% == == == == == == == == == == == == == == == == == == == == == == == == == ==
    == == == == == == == == == == == ==
    = % Section 17 : Software Implementation % == == == == == == == == == == ==
      == == == == == == == == == == == == == == == == == == == == == == == == ==
      == ==
    =

\subsection {
  Software Preliminary Implementation
}

\sectionauthors{BS, JB, CM}

This section covers the software architecture: operating system, SLAM, navigation, control, ML inference, and user interface. See Embedded Implementation for ESP32 firmware details.

\subsubsection{
  System Architecture Overview
}

\begin{table}[H]
\centering
\caption { Software Stack Layers }
\begin{tabular} {
  @{} llp {
    6cm
  }
  @{}
}
\toprule
\textbf{Layer} & \textbf{Platform} & \textbf { Responsibilities }
\\
\midrule Embedded &ESP32 - S3 &Motor PWM, sensor polling,
    hardware abstraction \ Edge &Raspberry Pi 5 & ROS 2, SLAM, Nav2,
    sensor fusion,
    local ML \ Remote(Optional) & Server &ML offloading, data storage, web UI \\
\bottomrule
\end {
  tabular
}
\end { table }

\begin{figure}[H]
\centering
\begin{tikzpicture}[box /.style = {rectangle, draw, rounded corners,
                                    minimum width = 2.2cm,
                                    minimum height = 0.8cm, align = center,
                                    font =\footnotesize, fill = white},
                     arrow /.style = {->, thick, >= stealth}] %
    ==
    = REMOTE LAYER == =
    \fill[blue !12, rounded corners](-5.2, 6.2) rectangle(5.2, 4.4);
\draw[dashed, blue !50, rounded corners](-5.2, 6.2) rectangle(5.2, 4.4);
\node[font =\small\bfseries, anchor = east] at(-5.4, 5.3){Remote};

\node[box](ui) at(-2.5, 5.3){Web UI\\(Vue + TS)};
\node[box](gateway) at(0.5, 5.3){Gateway\\(Kotlin)};
\node[box](ml_server) at(3.5, 5.3){ML Server\\(YOLO)};

% == = EDGE LAYER == =
    \fill[green !12, rounded corners](-5.2, 4.1) rectangle(5.2, 0.7);
\draw[dashed, green !50, rounded corners](-5.2, 4.1) rectangle(5.2, 0.7);
\node[font =\small\bfseries, anchor = east] at(-5.4, 2.4){Edge(Pi 5)};

\node[box](ros) at(-2.5, 3.3){ROS 2\\Humble};
\node[box](slam) at(0.5, 3.3){SLAM\\Algorithm};
\node[box](nav) at(3.5, 3.3){Nav2\\Stack};

\node[box](ml_local) at(-2.5, 1.5){ML Inference\\(Local)};
\node[box](fusion) at(0.5, 1.5){Sensor\\Fusion};
\node[box](control) at(3.5, 1.5){Motion\\Control};

% == = EMBEDDED LAYER == =
    \fill[orange !12, rounded corners](-5.2, 0.4) rectangle(5.2, -1.4);
\draw[dashed, orange !50, rounded corners](-5.2, 0.4) rectangle(5.2, -1.4);
\node[font =\small\bfseries, anchor = east] at(-5.4, -0.5){Embedded};

\node[box](esp) at(-1.5, -0.5){ESP32 - S3\\Firmware};
\node[box](motors) at(1.5, -0.5){Motor\\Drivers};
\node[box](sensors) at(4, -0.5){Sensors};

% == = CONNECTIONS == = % Remote horizontal
    \draw[arrow, <->](ui)--(gateway);
\draw[arrow, <->](gateway)--(ml_server);

% Remote to Edge
    \draw[arrow, <->](gateway)--(ros);

% Edge top row horizontal
    \draw[arrow](ros)--(slam);
\draw[arrow](slam)--(nav);

% Edge vertical
    \draw[arrow](nav)--(control);
\draw[arrow](fusion)--(slam);

% Edge bottom row horizontal
    \draw[arrow, <->](ml_local)--(fusion);

% Edge to Embedded
    \draw[arrow, <->](control)--(esp);

% Embedded horizontal
    \draw[arrow](esp)--(motors);
\draw[arrow](sensors.south)-- ++(0, -0.3) - | (esp.south);

\end { tikzpicture }
\caption { STAR Robot Software Architecture Overview }
\label {
fig:
  software - arch
}
\end { figure }

\subsubsection{Software Process Overview}

    The software system follows a modular 6 -
    layer architecture
        .Each layer operates independently while communicating through defined
            interfaces.The complete integrated process map is available in
                Appendix ~\ref {
app:
  software - process - full
} for
  full reference.

\paragraph{Layer 1 : User Interface}

      The UI layer provides real -
      time visualization and operator control
          through map display(800 $\times$600),
      camera feeds(640 $\times$480),
      and teleoperation inputs via keyboard / joystick at 50 ~Hz.

\input {
    sections / process_map_ui
  }

\paragraph{Layer 2: Navigation \& Control}

The navigation layer executes path planning via A* algorithm on the global costmap and local trajectory control using DWA for real-time obstacle avoidance.

\input{
  sections / process_map_navigation
}

\paragraph{Layer 3 : SLAM \& Mapping(Hector - SLAM)}

    Hector -
    SLAM processes LiDAR
    scans(SICK TiM561 at 15 ~Hz) to build occupancy grid maps while estimating
    robot pose through scan matching and loop closure detection.

\input {
  sections / process_map_slam
}

\paragraph{Layer 4 : Computer Vision Safety}

    The CV layer performs stereo depth
    estimation(640 $\times$480 @30 ~FPS) and
    YOLO - based person detection to enable safety -
        critical obstacle avoidance with emergency stop at distances below
        0.5 ~m.

\input {
  sections / process_map_cv
}

\paragraph{Layer 5 : Hardware \& Sensors}

This layer manages motor control via SPI,
    sensor polling(IMU via I2C, encoders, temperature via ADC),
    power system monitoring via SMBus,
    and publishes odometry to / odom and /
        tf.

\input {
  sections / process_map_hardware
}

\paragraph{Layer 6: ESP32-Firmware Demo (STAR Framework)}

The demo layer showcases STAR Framework capabilities with FreeRTOS tasks, DHT22 temperature/humidity sensors, HC-SR04 ultrasonic sensors, and PCA9685 PWM output for LED patterns.

\input{
  sections / process_map_demo
}

\subsubsection{Operating System}

Custom \textbf{Buildroot 2025.02 LTS} Linux for Raspberry Pi 5 \parencite{buildroot}:

\begin{table}[H]
\centering
\caption{
  Buildroot Image Configuration
}
\begin{tabular} {
  @{} ll @{}
}
\toprule
\textbf{Component} & \textbf { Details }
\\
\midrule Base image & 2 GB(vs.4 + GB Raspberry Pi OS) \ Boot time &
    15 seconds \ ROS 2 & Jazzy LTS with Nav2,
    SLAM Toolbox \ Vision &OpenCV 4 with DNN module \ ML Stack &Python 3.11,
    NumPy, SciPy, TFLite \ Networking &mDNS(star - robot.local), WiFi + Ethernet \\
\bottomrule
\end {
  tabular
}
\end{table}

        %
    == == == == == == == == == == == == == == == == == == == == == == == == ==
    == == == == == == == == == == == == ==
    = % SLAM ALGORITHM DECISION % == == == == == == == == == == == == == == ==
      == == == == == == == == == == == == == == == == == == == == == == ==
    =
\subsubsection{SLAM Algorithm Selection}

    Simultaneous Localization and
    Mapping(SLAM) is the core algorithmic capability enabling autonomous
        navigation \parencite{thrun_probabilistic_robotics}
            .We evaluated three prominent 2D LiDAR SLAM algorithms available in
                the ROS 2 ecosystem.

\paragraph{Candidate Algorithms}

    Three 2D LiDAR SLAM algorithms were evaluated :
\begin {
  itemize
}
\item \textbf { Hector SLAM }
\parencite{hector_slam} : Scan - matching without odometry;
low compute but no loop closure
    \item \textbf { SLAM Toolbox }
\parencite{slam_toolbox} : ROS 2 Nav2 default;
pose graph optimization with loop closure
    \item \textbf { Google Cartographer }
\parencite{cartographer} : Graph - based with submaps;
high accuracy but complex configuration
\end { itemize }

\paragraph { Decision Matrix }

\begin{table}[H]
\centering
\caption { SLAM Algorithm Decision Matrix }
\label {
tab:
  slam - decision
}
\begin{tabular} {
  @{} lcccc @{}
}
\toprule
\textbf{Criterion} & \textbf{Weight} & \textbf{Hector} & \textbf{
    SLAM Toolbox} & \textbf {
  Cartographer
}
\\
\midrule ROS 2 Nav2 Integration & 25\% &2 & \textbf{5} &
    4 \ Computational Efficiency & 20\% & \textbf{5} & 4 &
    2 \ Loop Closure Support & 15\% &1 & 4 & \textbf {
  5
}
\ Odometry Requirement & 15\% & \textbf{5} & 3 & 3 \ Documentation Quality &
    10\% &3 & \textbf{5} & 4 \ Configuration Complexity & 10\% & \textbf{5} &
    4 & 2 \ Active Development & 5\% &2 & \textbf{5} & 3 \\
\midrule
\textbf{Weighted Score} & 100\% &3.35 & \textbf{4.20} & 3.30 \\
\bottomrule
\end {
  tabular
}

\footnotesize {
Scoring:
  1 = Poor, 2 = Below Average, 3 = Average, 4 = Good, 5 = Excellent
}
\end { table }

\paragraph {
Selection:
  SLAM Toolbox
}

\textbf{SLAM Toolbox} selected for: native ROS 2 Nav2 integration, loop closure support, effective odometry fusion with our quadrature encoders, and balanced computational load on Pi 5. Hector SLAM remains as fallback if odometry proves unreliable.

\paragraph{SLAM Process Flow}

Figure~\ref{fig:slam-flowchart} illustrates the SLAM algorithm data flow, from sensor input through map generation.

\begin{figure}[H]
\centering
\resizebox{\textwidth}{
  !
}
{
  %
\begin{tikzpicture}[node distance = 0.8cm and 1.2cm,
                     block /.style = {rectangle, draw, rounded corners,
                                      minimum width = 1.8cm,
                                      minimum height = 0.7cm, align = center,
                                      font =\scriptsize},
                     decision /.style = {diamond, draw, minimum width = 1.2cm,
                                         minimum height = 0.6cm, align = center,
                                         font =\scriptsize, aspect = 2},
                     arrow /.style = {->, thick, >= stealth}] %
      Input sensors
    \node[block, fill = green !20](lidar){LiDAR\\Scan};
  \node[block, fill = green !20, below = 0.5cm of lidar](odom){Wheel\\Odometry};
  \node[block, fill = green !20, below = 0.5cm of odom](imu){IMU\\Data};

  % Preprocessing
    \node[block, fill = blue !20, right = 1.2cm of odom](preproc){
        Sensor\\Fusion};

  % Scan matching
    \node[block, fill = yellow !20, right = 1.2cm of preproc](scan){
        Scan\\Matching};

  % Loop closure decision
    \node[decision, fill = orange !20, right = 1.2cm of scan](loop){
        Loop\\Detected ? };

  % Graph optimization
    \node[block, fill = red !20, below = 0.8cm of loop](graph){
        Graph\\Optimization};

  % Map update
    \node[block, fill = purple !20, right = 1.2cm of loop](map){Update\\Map};

  % Output
    \node[block, fill = gray !20, right = 1cm of map](output){Occupancy\\Grid};

  % Connections
    \draw[arrow](lidar) - | (preproc);
  \draw[arrow](odom)--(preproc);
  \draw[arrow](imu) - | (preproc);
  \draw[arrow](preproc)--(scan);
  \draw[arrow](scan)--(loop);
  \draw[arrow](loop)-- node[above, font =\tiny]{No}(map);
  \draw[arrow](loop)-- node[right, font =\tiny]{Yes}(graph);
  \draw[arrow](graph) - | (map);
  \draw[arrow](map)--(output);

  % Feedback loop
    \draw[arrow, dashed](output.south)-- ++(0, -0.5) - | (scan.south);

  \end{tikzpicture} %
}
\caption { SLAM Toolbox Process Flow }
\label {
fig:
  slam - flowchart
}
\end{figure}

% =============================================================================
% LiDAR COMPARISON
% =============================================================================
\subsubsection{LiDAR Sensor Selection}

The STAR robot supports two LiDAR sensors, providing flexibility for different deployment scenarios. Both sensors have been validated with the ROS 2 software stack.

\paragraph{
  Sensor Comparison
}

\begin{table}[H]
\centering
\caption { LiDAR Sensor Comparison }
\label {
tab:
  lidar - comparison
}
\begin{tabular} {
  @{} lcc @{}
}
\toprule
\textbf{Specification} & \textbf{SICK TiM561} & \textbf { RPLiDAR C1 }
\\
\midrule Field of View & 270\textdegree & \textbf { 360\textdegree }
\ Range & 10 m & \textbf{12 m}(white),
    6 m(black) \ Scan Rate & \textbf{15 Hz} &
        10 Hz \ Angular Resolution & \textbf{0.33\textdegree}(810 points) &
        0.45\textdegree(800 points) \ Interface &Ethernet TCP
            / IP &UART(USB adapter) \ Accuracy & \textbf{$\pm$60 mm}(1 - 10 m) &
        $\pm$30 mm(0.15 - 0.5 m) \ IP Rating &IP67 &IP54 \ Operating Temp &
        -25\textdegree C to + 50\textdegree C &
        0\textdegree C to 40\textdegree C \ Power Consumption & 4 W & \textbf {
  2.5 W
}
\ Cost & \$3, 000 + (donated) & \textbf {\$60 }
\ ROS 2 Driver &sick\_scan\_xd &rplidar\_ros \\
\bottomrule
\end { tabular }
\end { table }

\paragraph{Selection Rationale}

    The STAR Robot uses a dual -
    sensor strategy.The \textbf{
        SICK TiM561} serves as the primary production sensor due to :
\begin {
  itemize
}
    \item Higher accuracy (0.33\textdegree{} angular resolution, $\pm$60mm distance accuracy)
    \item Industrial reliability (IP67 rating, automotive-grade)
    \item 15Hz scan rate for faster SLAM updates
    \item Donated by SICK Inc. (\$3,500+ value)
\end{itemize}

The \textbf{RPLiDAR C1} serves as a backup/development sensor, providing 360\textdegree{} FOV for software development and testing scenarios where the primary sensor is unavailable.

% =============================================================================
% ROS 2 ARCHITECTURE
% =============================================================================
\subsubsection{ROS 2 Architecture}

ROS 2 (Robot Operating System 2) provides the middleware framework for inter-process communication, sensor data distribution, and navigation algorithm coordination \parencite{ros2_documentation}. The STAR robot uses ROS 2 Jazzy, the latest LTS release.

\paragraph{Node Architecture}

The ROS 2 node graph defines how software components communicate through topics (publish-subscribe), services (request-response), and actions (long-running tasks with feedback).

\begin{figure}[H]
\centering
\resizebox{0.95\textwidth}{
      !
    }
    {
\begin{tikzpicture}[
    % Styles
    driver/.style={rectangle, draw, rounded corners, fill=green!20, minimum width=2.4cm, minimum height=1cm, align=center, font=\small},
    topic/.style={ellipse, draw, fill=blue!15, minimum width=2cm, minimum height=0.7cm, align=center, font=\scriptsize},
    process/.style={rectangle, draw, rounded corners, fill=yellow!30, minimum width=2.4cm, minimum height=1cm, align=center, font=\small},
    nav/.style={rectangle, draw, rounded corners, fill=orange!30, minimum width=2.4cm, minimum height=1cm, align=center, font=\small},
    ml/.style={rectangle, draw, rounded corners, fill=red!20, minimum width=2.4cm, minimum height=1cm, align=center, font=\small},
    arrow/.style={-Stealth, thick}
]
    % =====================================================================
    % COLUMN 1: Sensor Drivers (leftmost, shifted left for wire clearance)
    % =====================================================================
    \node[driver] (lidar) at (-3, 0) {LiDAR\\Driver};
\node[driver](imu) at(-3, -2.5){IMU\\Driver};
\node[driver](camera) at(-3, -5){Camera\\Driver};
\node[driver](esp) at(-3, -7.5){ESP32\\Bridge};

% == == == == == == == == == == == == == == == == == == == == == == == == == ==
    == == == == == == == ==
    = % COLUMN 2 : Input Topics % == == == == == == == == == == == == == == ==
      == == == == == == == == == == == == == == == == == == ==
    =
    \node[topic](scan) at(3.5, 0){ / scan};
\node[topic](imu_topic) at(3.5, -2.5){ / imu / data};
\node[topic](image) at(3.5, -5){ / camera / image};
\node[topic](odom) at(3.5, -7.5){ / odom};

% == == == == == == == == == == == == == == == == == == == == == == == == == ==
    == == == == == == == ==
    = % COLUMN 3 : Core Processing % == == == == == == == == == == == == == ==
      == == == == == == == == == == == == == == == == == == == ==
    =
    \node[process](amcl) at(7.5, 0){AMCL\\Localization};
\node[process](slam) at(7.5, -2.5){SLAM\\Toolbox};
\node[ml](yolo) at(7.5, -5){YOLO\\Detector};

% == == == == == == == == == == == == == == == == == == == == == == == == == ==
    == == == == == == == ==
    = % COLUMN 4 : Output Topics(between processing and nav) % == == == == == ==
      == == == == == == == == == == == == == == == == == == == == == == == == ==
      == == ==
    =
    \node[topic](map) at(10.5, -2.5){ / map};
\node[topic](detections) at(10.5, -5){ / detections};

% == == == == == == == == == == == == == == == == == == == == == == == == == ==
    == == == == == == == ==
    = % COLUMN 5 : Navigation % == == == == == == == == == == == == == == == ==
      == == == == == == == == == == == == == == == == == ==
    =
    \node[nav](nav_planner) at(13.5, 0){Nav2\\Planner};
\node[nav](nav_controller) at(13.5, -2.5){Nav2\\Controller};

% == == == == == == == == == == == == == == == == == == == == == == == == == ==
    == == == == == == == ==
    = % COLUMN 6 : Command output % == == == == == == == == == == == == == == ==
      == == == == == == == == == == == == == == == == == == ==
    =
    \node[topic](cmd_vel) at(16.5, -2.5){ / cmd\_vel};

% == == == == == == == == == == == == == == == == == == == == == == == == == ==
    == == == == == == == ==
    = % CONNECTIONS : Drivers to Topics % == == == == == == == == == == == == ==
      == == == == == == == == == == == == == == == == == == == == ==
    =
    \draw[arrow](lidar)--(scan);
\draw[arrow](imu)--(imu_topic);
\draw[arrow](camera)--(image);
\draw[arrow](esp)--(odom);

% == == == == == == == == == == == == == == == == == == == == == == == == == ==
    == == == == == == == ==
    = % CONNECTIONS : Topics to Processing Nodes % == == == == == == == == == ==
      == == == == == == == == == == == == == == == == == == == == == == == ==
    = % / scan to AMCL
    \draw[arrow](scan)--(amcl);

% / scan to SLAM
    \draw[arrow](scan) to[bend right = 20](slam);

% / imu / data to SLAM
    \draw[arrow](imu_topic)--(slam);

% / odom to SLAM
    \draw[arrow](odom) to[bend left = 60](slam);

% / odom to AMCL
    \draw[arrow](odom) to[bend left = 75](amcl);

% / camera / image to YOLO
    \draw[arrow](image)--(yolo);

% == == == == == == == == == == == == == == == == == == == == == == == == == ==
    == == == == == == == ==
    = % CONNECTIONS : Processing to Navigation % == == == == == == == == == ==
      == == == == == == == == == == == == == == == == == == == == == == == ==
    = % SLAM publishes / map
    \draw[arrow](slam)--(map);

% / map to Nav2 Planner
    \draw[arrow](map.north)-- ++(0, 0.8) - | (nav_planner.south);

% / map to Nav2 Controller
    \draw[arrow](map)--(nav_controller);

% AMCL to Nav2 Planner(robot pose)
    \draw[arrow](amcl)--(nav_planner);

% Nav2 Planner to Nav2 Controller
    \draw[arrow](nav_planner)--(nav_controller);

% Nav2 Controller to / cmd_vel
    \draw[arrow](nav_controller)--(cmd_vel);

% YOLO to / detections
    \draw[arrow](yolo)--(detections);

% == == == == == == == == == == == == == == == == == == == == == == == == == ==
    == == == == == == == ==
    = % FEEDBACK LOOP : / cmd_vel back to ESP32 Bridge % == == == == == == == ==
      == == == == == == == == == == == == == == == == == == == == == == == == ==
      ==
    =
    \draw[arrow, dashed, thick](cmd_vel.south)-- ++(0, -5.5) - | (esp.south);

% == == == == == == == == == == == == == == == == == == == == == == == == == ==
    == == == == == == == ==
    = % LEGEND % == == == == == == == == == == == == == == == == == == == == ==
      == == == == == == == == == == == == ==
    =
    \node[driver, minimum width = 1.5cm, minimum height = 0.5cm] at(1, -9.5){};
\node[font =\scriptsize, right] at(1.9, -9.5){Sensor Driver};

\node[topic, minimum width = 1.2cm, minimum height = 0.4cm] at(5, -9.5){};
\node[font =\scriptsize, right] at(5.7, -9.5){Topic};

\node[process, minimum width = 1.5cm, minimum height = 0.5cm] at(8.5, -9.5){};
\node[font =\scriptsize, right] at(9.4, -9.5){Processing};

\node[nav, minimum width = 1.5cm, minimum height = 0.5cm] at(12.5, -9.5){};
\node[font =\scriptsize, right] at(13.4, -9.5){Navigation};

\node[ml, minimum width = 1.5cm, minimum height = 0.5cm] at(16, -9.5){};
\node[font =\scriptsize, right] at(16.9, -9.5){ML};

\end { tikzpicture }
    }
    \caption {ROS 2 Node Graph for STAR Robot
    }
    \label {
    fig:
      ros2 - nodes
    }
    \end{figure}

            %
        == == == == == == == == == == == == == == == == == == == == == == == ==
        == == == == == == == == == == == == == ==
        = % NAV2 NAVIGATION STACK % == == == == == == == == == == == == == == ==
          == == == == == == == == == == == == == == == == == == == == == == ==
        =
\subsubsection{Navigation Stack(Nav2)}

    Nav2 is the ROS 2 navigation framework providing path planning,
        trajectory following,
        and behavior coordination \parencite{nav2_documentation, nav2_paper}.

\begin{table}[H]
\centering
\caption {
      Nav2 Component Configuration
    }
    \begin{tabular} {
      @{} llp {
        5.5cm
      }
      @{}
    }
    \toprule
\textbf{Component} & \textbf{Plugin} & \textbf { Function }
    \\
\midrule BT Navigator &NavigateToPose &Behavior tree
                coordination \ Planner Server &NavFn &Global A *path
                    planning \ Controller Server &DWA &Local trajectory +
            obstacle avoidance \ Costmap 2D &
        Layered &Static + LiDAR + inflation +
            ultrasonic \ Recovery Server &Spin,
        Backup &Unstick behaviors \\
\bottomrule
\end {
      tabular
    }
    \end { table }

    \paragraph{DWA Local Planner}

    The Dynamic Window
    Approach(DWA) is the local planner for real-time obstacle avoidance \parencite{dwa_fox}. At 20 Hz, DWA samples velocity commands within kinematic constraints, simulates trajectories, and selects the highest-scoring collision-free option.

\begin{table}[H]
\centering
\caption{
      DWA Configuration Parameters
    }
    \label {
    tab:
      dwa - params
    }
    \begin{tabular} {
      @{} llp {
        5cm
      }
      @{}
    }
    \toprule
\textbf{Parameter} & \textbf{Value} & \textbf { Rationale }
    \\
\midrule $v_{max} $ & 0.5 m / s &Safe indoor speed;
    25 cm stopping distance \ $\omega_{max} $ & 1.0 rad / s &
            180\textdegree ~in 3 s \ $\dot{v} _{max} $ & 0.5 m / s$ ^
        2 $ &Prevents wheel slip \ Robot footprint & 0.38 m $\times$ 0.38 m &
            1.25 ft square \ Simulation time & 2.0 s &Look - ahead horizon \\
\bottomrule
\end {
      tabular
    }
    \end { table }

    \textbf { Velocity Selection: }
    0.5 m / s max ensures adequate sensor
                response(5 cm travel between 10 Hz LiDAR scans) and
        stopping distance($d = v ^ 2 / 2a = 0.25 $ m at 0.5 m / s$ ^ 2 $ decel)
            .

\paragraph{DWA Algorithm Flow}

        The DWA controller samples velocities,
        simulates trajectories,
        and selects the optimal path at 20 Hz :

\begin{figure}[H]
\centering
\begin{tikzpicture}[node distance = 0.6cm and 1.2cm,
                     block /.style = {rectangle, draw, rounded corners,
                                      minimum width = 1.8cm,
                                      minimum height = 0.6cm, align = center,
                                      font =\tiny},
                     decision /.style = {diamond, draw, minimum width = 1.2cm,
                                         minimum height = 0.6cm, align = center,
                                         font =\tiny, aspect = 2},
                     arrow /.style = {->, thick, >= stealth}] %
            Input
    \node[block, fill = green !20](goal){Goal\\Pose};
    \node[block, fill = green !20, below = 0.4cm of goal](costmap){
        Costmap\\(Obstacles)};
    \node[block, fill = green !20, below = 0.4cm of costmap](odom){
        Current\\Velocity};

    % Sample
    \node[block, fill = blue !20, right = 1.5cm of costmap](sample){
          Sample\\Velocities};

    % Simulate
    \node[block, fill = yellow !20, right = 1.2cm of sample](sim){
          Simulate\\Trajectories};

    % Score
    \node[block, fill = orange !20, right = 1.2cm of sim](score){Score\\Paths};

    % Check collision
    \node[decision, fill = red !20, right = 1.2cm of score](collision){
          Collision\\Free ? };

    % Select best
    \node[block, fill = purple !20, above = 0.6cm of collision](best){
          Select\\Best};

    % Output
    \node[block, fill = gray !20, right = 1.2cm of collision](output){
          cmd\_vel\\Output};

    % Connections
    \draw[arrow](goal) - | (sample);
    \draw[arrow](costmap)--(sample);
    \draw[arrow](odom) - | (sample);
    \draw[arrow](sample)--(sim);
    \draw[arrow](sim)--(score);
    \draw[arrow](score)--(collision);
    \draw[arrow](collision)-- node[above, font =\tiny]{Yes}(output);
    \draw[arrow](collision)-- node[left, font =\tiny]{No}(best);
    \draw[arrow](best) - | (sim);

    \end { tikzpicture }
    \caption { DWA Local Planner Algorithm Flow }
    \label {
    fig:
      dwa - flow
    }
    \end { figure }

\paragraph{Ultrasonic Integration}

Seven HC-SR04 sensors provide short-range detection for DWA, particularly for obstacles below the LiDAR plane (table legs, pets).

\begin{table}[H]
\centering
\caption{
  HC - SR04 Range Utilization
}
\begin{tabular} {
  @{} lll @{}
}
\toprule
\textbf{Range} & \textbf{Accuracy} & \textbf { Action }
\\
\midrule 2 --30 cm &$\pm$3 mm &Emergency stop \ 30 --150 cm &$\pm$1 cm &
        DWA costmap integration \ $ > $150 cm &$\pm$3 + cm &Ignored(unreliable) \\
\bottomrule
\end {
  tabular
}
\end{table}

        %
    == == == == == == == == == == == == == == == == == == == == == == == == ==
    == == == == == == == == == == == == ==
    = % CONTROL SYSTEMS % == == == == == == == == == == == == == == == == == ==
      == == == == == == == == == == == == == == == == == == == ==
    =
\subsubsection{Control Systems}

    The STAR robot implements a hierarchical control architecture with cascaded
        control loops operating at different frequencies.

\paragraph{Control Hierarchy}

    The control system uses established robotics control theory to implement
        velocity control and
    navigation.The following equations,
    derived from standard references \parencite{astrom_pid, siciliano_robotics},
    will be used to model the control system during implementation.

\begin{figure}[H]
\centering
\resizebox{0.95\textwidth} {
  !
}
{
  %
\begin{tikzpicture} % Input from navigation goal
    \sbEntree{E}

      % Main control chain
    \sbBloc[3] {
    nav
  }
  { Nav2 Planner }
  { E }
  \sbBloc[3] { dwa }
  { DWA Controller }
  { nav }
  \sbComp[3] { sum }
  { dwa }
  \sbBloc[3] { pid }
  { Wheel PID }
  { sum }
  \sbBloc[3] { motor }
  { Motor Driver }
  { pid }
  \sbSortie[3] { S }
  {motor}

      % Connect the chain
    \sbRelier{E} {
    nav
  }
  \sbRelier{nav} { dwa }
  \sbRelier{dwa} { sum }
  \sbRelier{sum} { pid }
  \sbRelier{pid} { motor }
  \sbRelier{motor} {S}

      % Labels on connections
    \sbNomLien[0.5] {
    E
  }
  { Goal }
  \sbNomLien[0.5] { S }
  {Motion}

      % Encoder feedback loop
    \sbDecaleNoeudy[-4] {
    pid
  }
  { encfb }
  \sbBlocr{enc} { Encoder }
  { encfb }
  \sbRelieryx{motor - S} { enc }
  \sbRelierxy{enc} {sum}

      % Rate labels above blocks
    \node[font =\scriptsize, above = 0.2cm of nav]{10 Hz};
  \node[font =\scriptsize, above = 0.2cm of dwa]{20 Hz};
  \node[font =\scriptsize, above = 0.2cm of pid]{100 Hz};
  \end{tikzpicture} %
}
\caption { Hierarchical Control Architecture }
\label {
fig:
  control - hierarchy
}
\end { figure }

\paragraph{Velocity Control(PID)}

Each wheel will be controlled by an independent PID velocity controller running
    on the ESP32 at 100 Hz \parencite{astrom_pid} :

\begin{equation} u(t) = K_p e(t) + K_i \int_0
                         ^ t e(\tau) d\tau + K_d \frac{de(t)} {
  dt
}
\end{equation}

where $u(t)
$ is the PWM duty cycle,
    $e(t) = v_{cmd} - v_{actual} $ is the velocity error, and $K_p$, $K_i$,
    $K_d$ are tunable gains.

\paragraph{Differential Drive Kinematics}

        The 4 -
        wheel differential drive configuration follows the standard 2 -
        wheel kinematic model \parencite{siciliano_robotics} :

\textbf {
  Forward Kinematics:
}
\begin{align} v &= \frac{v_R + v_L} {2}, \quad \omega = \frac{v_R - v_L} { L }
\end { align }

\textbf { Inverse Kinematics: }
\begin{align} v_R &= v + \frac{\omega L} {2}, \quad v_L = v - \frac{\omega L} {
  2
}
\end{align}

where $v$ is linear velocity,
    $\omega$ is angular velocity, $v_R$ / $v_L$ are wheel velocities,
    and $L$ is the effective wheel base.

\paragraph{Odometry Calculation}

        Wheel odometry integrates encoder counts to estimate robot pose
        $(x, y, \theta) $ :
\begin {
  align
}
\Delta s &= \frac{\Delta s_R + \Delta s_L} {2}, \quad \Delta \theta = \frac{\Delta s_R - \Delta
                                                                                s_L} {
  L
}
\end { align }
\begin{align}
x_{k+1} &= x_k + \Delta s \cdot \cos(\theta_k + \Delta\theta/2), \quad y_{k+1} = y_k + \Delta s \cdot \sin(\theta_k + \Delta\theta/2)
\end{align}

% =============================================================================
% MACHINE LEARNING - PERSON DETECTION
% =============================================================================
\subsubsection{Machine Learning: Person Detection}

The STAR robot incorporates real-time person detection to support first responder applications. Detected persons are marked on the generated map, providing situational awareness for emergency response teams.

\paragraph{
  Use Case Requirements
}

\begin { itemize }
    \item \textbf{Primary Function}: Detect and localize people within camera field of view
    \item \textbf{Output}: Position markers on the occupancy grid map indicating person locations
    \item \textbf{Safety Consideration}: Robot should slow or stop when persons are detected in the navigation path
    \item \textbf{Data Logging}: Detection timestamps and estimated positions stored for post-mission analysis
\end{
      itemize
    }

\paragraph{YOLO Model Selection}

YOLO (You Only Look Once) is the selected object detection framework due to its real-time performance characteristics \parencite{yolov8_ultralytics}. We evaluated multiple YOLO variants for deployment on the Raspberry Pi 5.

\begin{table}[H]
\centering
\caption{YOLO Model Comparison for Raspberry Pi 5
}
\label {
tab:
  yolo - comparison
}
\begin{tabular} {
  @{} lcccc @{}
}
\toprule
\textbf{Model} & \textbf{Parameters} & \textbf { mAP @50 }
& \textbf{Pi 5 FPS} & \textbf { Memory }
\\
\midrule YOLOv8n & 3.2M & 37.3\% &8 --12 & 6.3 MB \ YOLOv8s & 11.2M &
    44.9\% &3 --5 & 22.5 MB \ YOLOv8m & 25.9M & 50.2\% &1 --2 &
    52.0 MB \ YOLOv8l & 43.7M & 52.9\% &$ < $1 & 87.7 MB \ YOLOv8n(INT8) &
    3.2M & 35.5\% &15 --20 & 3.2 MB \\
\bottomrule
\end {
  tabular
}

\footnotesize {FPS estimates for 640$\times$480 input on Pi 5 CPU. INT8 quantization uses TensorFlow Lite. mAP@50 = mean Average Precision at 50\% IoU threshold on COCO dataset.
}
\end { table }

\paragraph { Decision Matrix }

\begin{table}[H]
\centering
\caption { YOLO Model Decision Matrix }
\label {
tab:
  yolo - decision
}
\begin{tabular} {
  @{} lccccc @{}
}
\toprule
\textbf{Criterion} & \textbf{Weight} & \textbf{v8n} & \textbf{v8s} & \textbf{
    v8m} & \textbf {
  v8n INT8
}
\\
\midrule Real - time Performance($\geq$5 FPS) & 30\% &4 & 2 & 1 & \textbf {
  5
}
\ Detection Accuracy & 25\% &3 & 4 & \textbf{5} & 3 \ Memory Footprint &
    20\% &4 & 3 & 2 & \textbf {
  5
} \\
CPU Headroom (for SLAM) & 15\% & 3 & 2 & 1 & \textbf{5}
\ Implementation Complexity & 10\% & \textbf{5} & 5 & 5 & 3 \\
\midrule
\textbf{Weighted Score} & 100\% &3.60 & 2.95 & 2.40 & \textbf {
  4.30
}
\\
\bottomrule
\end { tabular }
\end { table }

\paragraph{Selection: YOLOv8n INT8}

YOLOv8n with INT8 quantization selected for Pi 5: 15--20 FPS, 3.2 MB model, minimal CPU impact (30--50\%), 1.8\% mAP reduction acceptable for person detection. Uses TensorFlow Lite runtime \parencite{tflite_quantization}.

\paragraph{
  Local vs.Remote Inference
}

\begin{table}[H]
\centering
\caption { Local vs.Remote Inference Trade - offs }
\label {
tab:
  inference - location
}
\begin{tabular} {
  @{} lcc @{}
}
\toprule
\textbf{Factor} & \textbf{Pi 5(Local)} & \textbf { Server(Remote) }
\\
\midrule Latency & 50 --100 ms & 100 --300 ms(network dependent) \ FPS &
    15 --20(INT8) &
    60 + (FP16) \ Model Size &YOLOv8n only &YOLOv8m / l possible \ Accuracy &
    35.5\% mAP & 50 + \% mAP \ Network Dependency &None &Required \ CPU Impact &
    30 --50\% &Minimal \\
\bottomrule
\end {
  tabular
}
\end{table}

Runtime switching between local/remote based on network availability and CPU load.

% =============================================================================
% REMOTE SERVER ARCHITECTURE
% =============================================================================
\subsubsection{Remote Server Architecture}

Optional server for ML offloading, data persistence, and UI hosting.

\begin{table}[H]
\centering
\caption{
  Server Component Stack
}
\begin{tabular} {
  @{} lllp {
    5cm
  }
  @{}
}
\toprule
\textbf{Component} & \textbf{Technology} & \textbf{Protocol} & \textbf {
  Function
}
\\
\midrule Gateway &Kotlin / Ktor &WebSocket &ROS 2 bridge, session mgmt,
    REST API,
    JWT auth \ ML Server &PyTorch / CUDA &gRPC &YOLOv8m /
            l inference on RTX 5070 \ Database &PostgreSQL +
        PostGIS &SQL &Maps,
    detections, trajectories,
    metadata \ Web UI &Vue 3 + TypeScript &HTTP / WS &Real -
        time monitoring and control \\
\bottomrule
\end {
  tabular
}
\end { table }

\paragraph { Communication Protocol }

\begin{table}[H]
\centering
\caption { Robot - Server Data Streams }
\begin{tabular} {
  @{} lll @{}
}
\toprule
\textbf{Stream} & \textbf{Rate} & \textbf { Content }
\\
\midrule Telemetry & 10 Hz &Pose, velocity,
    sensor status \ Map Updates &On change &
            Incremental occupancy grid deltas \ Camera Frames &
        5 FPS &JPEG -
            compressed(when remote ML enabled) \ Commands &On demand &Velocity,
    goals, mode switches \ Heartbeat & 1 Hz &Keep - alive with auto - reconnect \\
\bottomrule
\end {
  tabular
}
\end { table }

4G cellular backup available when WiFi unavailable (+100--500 ms latency).

% =============================================================================
% USER INTERFACE
% =============================================================================
\subsubsection{User Interface}

Web-based UI for real-time monitoring, control, and visualization. Built with Vue 3, TypeScript, Pinia, Tailwind CSS, and WebSocket streaming \parencite{vue3_documentation, typescript_documentation}.

\paragraph{Web UI Architecture}

The STAR Web UI follows a reactive single-page application architecture. Figure~\ref{fig:webui-flowchart} illustrates the data flow from robot sensors through the UI rendering pipeline.

\begin{figure}[H]
\centering
\resizebox{\textwidth}{
  !
}
{
  %
\begin{tikzpicture}[node distance = 0.8cm and 1.5cm,
                     block /.style = {rectangle, draw, rounded corners,
                                      minimum width = 2cm,
                                      minimum height = 0.7cm, align = center,
                                      font =\scriptsize},
                     store /.style = {rectangle, draw, rounded corners,
                                      minimum width = 2cm,
                                      minimum height = 0.7cm, align = center,
                                      font =\scriptsize, fill = purple !20},
                     component /.style =
                         {rectangle, draw, rounded corners, minimum width = 2cm,
                          minimum height = 0.7cm, align = center,
                          font =\scriptsize, fill = blue !20},
                     service /.style = {rectangle, draw, rounded corners,
                                        minimum width = 2cm,
                                        minimum height = 0.7cm, align = center,
                                        font =\scriptsize, fill = green !20},
                     arrow /.style = {->, thick, >= stealth}] %
      Robot Layer
    \fill[orange !12, rounded corners](-6.5, 4.5) rectangle(-2.5, 2.5);
  \draw[dashed, orange !50, rounded corners](-6.5, 4.5) rectangle(-2.5, 2.5);
  \node[font =\scriptsize\bfseries] at(-4.5, 4.2){Robot(Pi 5)};

  \node[block, fill = orange !20](ros) at(-4.5, 3.3){ROS 2\\Topics};

  % Server Layer
    \fill[green !12, rounded corners](-1.5, 4.5) rectangle(2.5, 2.5);
  \draw[dashed, green !50, rounded corners](-1.5, 4.5) rectangle(2.5, 2.5);
  \node[font =\scriptsize\bfseries] at(0.5, 4.2){Gateway Server};

  \node[service](gateway) at(0.5, 3.3){Kotlin / Ktor\\WebSocket};

  % Browser Layer
    \fill[blue !12, rounded corners](3.5, 4.5) rectangle(9.5, -1.5);
  \draw[dashed, blue !50, rounded corners](3.5, 4.5) rectangle(9.5, -1.5);
  \node[font =\scriptsize\bfseries] at(6.5, 4.2){Browser(Vue 3 + TypeScript)};

  % State Management
    \node[store](pinia) at(5, 3.3){Pinia\\Store};

  % Composables
    \node[service](useRobot) at(8, 3.3){useRobot()\\Composable};

  % UI Components Row 1
    \node[component](dashboard) at(4.5, 1.5){Dashboard\\Component};
  \node[component](mapping) at(6.5, 1.5){Mapping\\Component};
  \node[component](sensors) at(8.5, 1.5){Sensor Data\\Component};

  % UI Components Row 2
    \node[component](control) at(4.5, 0){Control\\Component};
  \node[component](analytics) at(6.5, 0){Analytics\\Component};
  \node[component](settings) at(8.5, 0){Settings\\Component};

  % Connections
    \draw[arrow](ros)--(gateway);
  \draw[arrow](gateway)--(pinia);
  \draw[arrow, <->](pinia)--(useRobot);

  % Pinia to components
    \draw[arrow](pinia)--(dashboard);
  \draw[arrow](pinia)--(mapping);
  \draw[arrow](useRobot)--(sensors);

  \draw[arrow](pinia)--(control);
  \draw[arrow](pinia)--(analytics);
  \draw[arrow](useRobot)--(settings);

  \end{tikzpicture} %
}
\caption { Web UI Data Flow Architecture }
\label {
fig:
  webui - flowchart
}
\end { figure }

\paragraph{Dashboard and Sensor Views}

    The Dashboard
    page(Figure ~\ref{fig : webui - dashboard}) provides a real
    - time mapping visualization with concentric range rings(2 --10 m),
    robot position indicator,
    and live LiDAR scan overlay.Status cards display battery
        level(with estimated runtime),
    WiFi signal strength, connection status,
    and current operating mode.The IMU Orientation panel shows pitch, roll,
    yaw angles and velocity with a 3D model preview indicating robot attitude.

\begin{figure}[H]
\centering
\includegraphics[width = 0.95\textwidth] {
  figures / webui_dashboard.png
}
\caption {
  STAR Web UI Dashboard : Real - time mapping view with IMU orientation display
}
\label {
fig:
  webui - dashboard
}
\end{figure}

The Sensor Data page (Figure~\ref{fig:webui-sensor-data}) displays all telemetry streams: GPS coordinates (latitude, longitude, altitude), internal temperature with thermal warning indicators, CPU/RAM utilization gauges, and motor load percentage. The Raw Data Log provides timestamped system events for debugging and mission review.

\begin{figure}[H]
\centering
\includegraphics[width=0.95\textwidth]{
  figures / webui_sensor_data.png
}
\caption {
  STAR Web UI Sensor Data : Telemetry dashboard with GPS,
                            thermal,
                            and system metrics
}
\label {
fig:
  webui - sensor - data
}
\end { figure }

\begin{table}[H]
\centering
\caption { Web UI Pages }
\label {
tab:
  ui - pages
}
\begin{tabular} {
  @{} lp {
    8cm
  }
  @{}
}
\toprule
\textbf{Page} & \textbf { Features }
\\
\midrule Dashboard &Video feed, system status, battery,
    connection indicator \ Mapping &Live occupancy grid, robot pose,
    LiDAR overlay, person markers, path history \ Control &Virtual joystick,
    WASD keys,
    speed presets(0.1 / 0.3 / 0.5 m / s), e - stop \ Sensors &LiDAR scan,
    IMU orientation, ultrasonic readings \ Analytics &Mission stats, coverage \%
    , detection counts \ Settings &Robot config, network,
    map export(PGM, PNG, PDF, GeoJSON) \\
\bottomrule
\end {
  tabular
}
\end{table}

Responsive design
    : desktop(full dashboard),
      tablet(field - optimized),
      mobile(control - focused).

              %
          == == == == == == == == == == == == == == == == == == == == == == ==
          == == == == == == == == == == == == == == ==
      = % SUMMARY % == == == == == == == == == == == == == == == == == == == ==
        == == == == == == == == == == == == == == == == == ==
      =
\subsubsection {
  Software Implementation Summary
}

\begin{table}[H]
\centering
\caption { Software Stack Summary }
\begin{tabular} {
  @{} ll @{}
}
\toprule
\textbf{Subsystem} & \textbf { Implementation }
\\
\midrule OS &Custom Buildroot Linux(
    ROS 2 Jazzy) \ SLAM &SLAM Toolbox with loop closure \ LiDAR &RPLiDAR
        C1(prototype) /
        SICK TiM561(industrial) \ Navigation &Nav2
    + DWA +
    layered costmaps \ Control &Hierarchical PID at
    100 Hz \ ML &YOLOv8n INT8 local;
GPU offload optional \ Server &Kotlin gateway + Vue UI \\
\bottomrule
\end {
  tabular
}
\end { table }

\textbf{Demonstration Goals : } Autonomous mapping,
    DWA obstacle avoidance with sensor fusion, real - time person detection,
    web - based teleoperation.
